% Section 10.8, from lectures/L10/appendix/b_can_def_package.md.

\section{The \code{can_def} package}
\label{c10:sec:candef}

\subsection*{Why use a VHDL package?}
\begin{itemize}
  \item Part~\ref{part:vhdl} never needed one: every example there was self-contained, and where an
    exercise reused a module you had already written \textemdash{} \code{reset_sync},
    \code{button_sync}, \code{timer} \textemdash{} you copied the file into the directory that
    needed it.
  \item This project is different. \code{TICKS_PER_BIT}, \code{SAMPLE_TICK}, \code{ID_WIDTH},
    \code{CRC_WIDTH} and more are each needed by \emph{several} of the blocks built from
    \chapref{c12:ch} onwards, and they have to agree exactly, everywhere, or the design breaks
    silently. That much shared information is worth defining once.
  \item Every constant in it comes from something this chapter has just taught: the field widths are
    the ones you lay out by hand in the exercises, and \code{SAMPLE_TICK} is the 70~\% point
    \secref{c10:sec:bittiming} reasoned about. This package is where those numbers stop being prose
    and become the single source of the definitions every module uses.
  \item It also changes the layout of the \emph{files}. Every module here is part of \textbf{one}
    design, so all of them live in \file{controller/} and read the single \file{can_def.vhd} beside
    them. Nothing is duplicated, and there is exactly one file to change when a constant changes.
\end{itemize}

\subsection*{Reading \code{can_def.vhd}}
Live-coded together during the session, and written \textbf{once}, to \file{controller/can_def.vhd}.
This is the first VHDL the book writes. Every module you write from \chapref{c11:ch} onwards lands in
the same directory, and all of them except \code{meta_prev} read this one file.

\textbf{The bit timing constants}, derived rather than hand-picked, so that they stay consistent with
each other:

\begin{vhdlcode}
constant CLOCK_FREQ_HZ: natural := 50_000_000;                  -- System clock frequency: 50 MHz.
constant BIT_RATE_HZ  : natural := 1_000_000;                   -- CAN bit rate: 1 Mbit/s.
constant TICKS_PER_BIT: natural := CLOCK_FREQ_HZ / BIT_RATE_HZ; -- Clock ticks per CAN bit: 50.
constant SAMPLE_TICK  : natural := (TICKS_PER_BIT * 7) / 10;    -- Sample point: 35 (70%).
\end{vhdlcode}

\code{TICKS_PER_BIT} has to divide evenly for this simple design to hit an exact bit rate:
$50\,000\,000 / 1\,000\,000$ is exactly 50. \code{SAMPLE_TICK} is deliberately computed as a
\emph{percentage} of \code{TICKS_PER_BIT}, not as a hardcoded tick number, so that a later change of
bit rate does not silently leave the sample point wrong. This is
\secref{c10:sec:bittiming}'s 70~\% sample point, turned into a number \code{bit_timer} can count to
in \chapref{c12:ch}.

\textbf{The frame field widths}, straight from the frame diagram in \secref{c10:sec:frameformat}:

\begin{vhdlcode}
constant ID_WIDTH  : natural := 11;          -- Standard CAN identifier width in bits.
constant DLC_WIDTH : natural := 4;           -- Data Length Code field width in bits.
constant CRC_WIDTH : natural := 15;          -- CAN CRC-15 sequence width in bits.
constant DLC_MAX   : natural := 8;           -- Maximum payload length in bytes.
constant DATA_WIDTH: natural := DLC_MAX * 8; -- Maximum data-field width in bits.
\end{vhdlcode}

\textbf{The bit stuffing rule's run length}, from \secref{c10:sec:stuffing}'s rule about five
identical bits:

\begin{vhdlcode}
constant MAX_RUN: natural := 5; -- Maximum run of identical bits before a stuff bit is inserted.
\end{vhdlcode}

Both shift registers count runs against this single number (\chapref{c14:ch}, \ref{c15:ch}): the
transmitter to decide when a stuff bit must be \emph{inserted}, the receiver to predict where one
must \emph{appear}. That both ends of the link agree on it exactly is this whole package's argument
in miniature \textemdash{} one definition, several modules, no way to drift apart.

\textbf{Common subtypes}, so that every module names identically shaped signals the same way:

\begin{vhdlcode}
subtype byte_t is std_logic_vector(7 downto 0);            -- 8-bit byte.
subtype id_t   is std_logic_vector(ID_WIDTH-1 downto 0);   -- Standard CAN identifier.
subtype dlc_t  is std_logic_vector(DLC_WIDTH-1 downto 0);  -- Data Length Code field.
subtype crc_t  is std_logic_vector(CRC_WIDTH-1 downto 0);  -- CAN CRC-15 sequence.
subtype data_t is std_logic_vector(DATA_WIDTH-1 downto 0); -- CAN data field.
\end{vhdlcode}

Every module except \code{meta_prev} starts with \code{use work.can_def.all;} and uses these names
directly: \code{bit_timer}'s counter range, \code{crc15}'s register width, and
\code{tx_shift_reg}/\code{rx_shift_reg}'s byte ports. All of it goes back to this one file.

\textbf{The CRC-15 polynomial}, as a 15-bit-wide constant. \Chapref{c13:ch}'s \code{crc15} uses it
directly; the mathematics is covered there, not here. For now treat it as a fixed standard constant
that every CAN controller uses:

\begin{vhdlcode}
constant CRC_POLY: crc_t := "100010110011001";
\end{vhdlcode}

It corresponds to $x^{15} + x^{14} + x^{10} + x^{8} + x^{7} + x^{4} + x^{3} + 1$.

\subsection*{What comes next}
\Chapref{c11:ch} maps this chapter's protocol onto a block architecture and writes the controller's
top level, \file{can_controller.vhd}, whose ports are typed from this package (\code{id_t},
\code{dlc_t}, \code{data_t}). It also introduces the simulation flow and leaves that top level's
architecture empty. \Chapref{c12:ch} builds the first two blocks to go into it,
\file{meta_prev.vhd} and \file{bit_timer.vhd}.
